\chapter{Cơ hội và quyết định}
\markboth{Cơ hội và quyết định}{Opportunity and Choice}

\section{Một người bỏ qua việc học mới vì nghĩ mai vẫn còn kịp.}
\begin{truyen}{Một người bỏ qua việc học mới vì nghĩ mai vẫn còn kịp.}{Opportunity rarely shouts twice.}
\chuhoa{C}{ô được mời tham gia một khóa học tốt nhưng ngại bắt đầu lại. Cô trì hoãn tới khi chương trình đóng tuyển. Lúc đó cô mới hiểu cơ hội không phải lúc nào cũng quay lại gõ cửa với cùng một sự rộng lượng.}
\textit{(She was invited to join a valuable course but hesitated to begin again. She delayed until registration closed. Then she learned that opportunity does not always return and knock with the same generosity.)}
\end{truyen}
\begin{baihoc}
Có những cánh cửa không đóng sầm; chúng chỉ khép dần cho đến khi ta nhận ra đã muộn.
\textit{(Some doors do not slam shut; they simply close little by little until it is too late.)}
\end{baihoc}
\begin{ghinhoanh}
\textbf{Short English to remember:}

Opportunity rarely shouts twice.

Move while the door is still open.
\end{ghinhoanh}
\begin{tuhocnhanh}
1. Cơ hội nào của em đang bị trì hoãn vì sợ bắt đầu? \textit{(What opportunity are you delaying because you fear starting?)}

2. Nếu mất nó, em có hối tiếc không? \textit{(Would you regret it if you lost it?)}
\end{tuhocnhanh}
\ngancach

\section{Một quyết định năm phút làm đổi hướng cả vài năm.}
\begin{truyen}{Một quyết định năm phút làm đổi hướng cả vài năm.}{Small choices cast long shadows.}
\chuhoa{M}{ột lời đồng ý bốc đồng kéo anh vào công việc không phù hợp, môi trường không lành mạnh, và chuỗi ngày mất phương hướng. Sau này anh mới thấy đời người nhiều khi không rẽ ở những quyết định rất lớn, mà ở những phút chọn đại cho xong.}
\textit{(One impulsive yes pulled him into the wrong job, an unhealthy environment, and years of confusion. Later he saw that life often does not turn at huge decisions, but at careless small ones.)}
\end{truyen}
\begin{baihoc}
Những lựa chọn nhỏ lặp lại thường tạo ra số phận rõ hơn những lời tuyên bố lớn.
\textit{(Repeated small choices shape destiny more clearly than big declarations.)}
\end{baihoc}
\begin{ghinhoanh}
\textbf{Short English to remember:}

Small choices cast long shadows.

Choose carefully when it looks small.
\end{ghinhoanh}
\begin{tuhocnhanh}
1. Gần đây em có quyết định nhỏ nào nhưng hậu quả có thể dài lâu? \textit{(Have you made a small recent decision with long-term consequences?)}

2. Em có đang chọn đại vì mệt hay vì thật sự thấy đúng? \textit{(Are you choosing carelessly because you are tired, or because it truly feels right?)}
\end{tuhocnhanh}
\ngancach

\section{Người trẻ từ chối cơ hội vì sợ mình chưa đủ giỏi.}
\begin{truyen}{Người trẻ từ chối cơ hội vì sợ mình chưa đủ giỏi.}{Sometimes readiness grows after yes.}
\chuhoa{C}{ậu đợi tới lúc hoàn hảo mới dám bước vào việc lớn. Nhưng cuộc đời không thường trao cơ hội cho bản hoàn hảo; nó trao cho người chịu lớn lên trong quá trình làm.}
\textit{(He waited until perfection before entering something bigger. But life rarely gives opportunities to finished versions of us; it gives them to people willing to grow during the process.)}
\end{truyen}
\begin{baihoc}
Nhiều người bỏ lỡ không phải vì quá yếu, mà vì đòi mình phải sẵn sàng tuyệt đối mới được bắt đầu.
\textit{(Many people miss out not because they are too weak, but because they demand perfect readiness before they begin.)}
\end{baihoc}
\begin{ghinhoanh}
\textbf{Short English to remember:}

Sometimes readiness grows after yes.

Do not wait to feel perfect.
\end{ghinhoanh}
\begin{tuhocnhanh}
1. Em đang từ chối điều gì chỉ vì thấy mình chưa đủ? \textit{(What are you refusing only because you feel not enough?)}

2. Em có thể học trong lúc làm được không? \textit{(Can you learn while doing it?)}
\end{tuhocnhanh}
\ngancach

\section{Một người luôn đổi ý vì sợ quyết rồi sai.}
\begin{truyen}{Một người luôn đổi ý vì sợ quyết rồi sai.}{Indecision is also a decision.}
\chuhoa{V}{ì sợ chọn nhầm, cô chần chừ ở mọi ngã rẽ. Cuối cùng, chính sự không chọn khiến thời cơ qua mất. Cô mới hiểu đứng yên quá lâu cũng là một cách để tự đánh mất đường đi.}
\textit{(Because she feared choosing wrong, she hesitated at every turning point. In the end, the opportunity passed precisely because she never chose. She learned that standing still too long is also a way of losing the path.)}
\end{truyen}
\begin{baihoc}
Không quyết định không giúp ta an toàn hơn nhiều; nhiều khi nó chỉ khiến ta mất quyền chủ động.
\textit{(Not deciding does not make us much safer; often it only makes us lose our agency.)}
\end{baihoc}
\begin{ghinhoanh}
\textbf{Short English to remember:}

Indecision is also a decision.

Standing still has a cost.
\end{ghinhoanh}
\begin{tuhocnhanh}
1. Em đang chần chừ ở ngã rẽ nào? \textit{(At which turning point are you hesitating?)}

2. Cái giá của việc không chọn là gì? \textit{(What is the cost of not choosing?)}
\end{tuhocnhanh}
\ngancach

\section{Người ta hiểu muộn rằng cơ hội thường mặc áo công việc khó.}
\begin{truyen}{Người ta hiểu muộn rằng cơ hội thường mặc áo công việc khó.}{Opportunity often looks like hard work.}
\chuhoa{N}{hiều việc ban đầu nhìn không hào nhoáng: phải học thêm, phải bắt đầu nhỏ, phải chịu áp lực. Vì thế người ta bỏ qua. Rồi về sau mới hiểu chính những thứ khó nhằn ấy lại là lối mở cho đời mình.}
\textit{(Many opportunities do not look glamorous at first: they require extra learning, humble beginnings, and pressure. So people pass them by, only to realize later that these hard things were actually openings.)}
\end{truyen}
\begin{baihoc}
Cơ hội hiếm khi đến với bảng hiệu dễ chịu; nó thường đến dưới dạng trách nhiệm khó nhằn.
\textit{(Opportunity rarely arrives with a comfortable signboard; it often comes as demanding responsibility.)}
\end{baihoc}
\begin{ghinhoanh}
\textbf{Short English to remember:}

Opportunity often looks like hard work.

Do not reject the rough beginning.
\end{ghinhoanh}
\begin{tuhocnhanh}
1. Em có đang bỏ qua một việc khó nhưng có giá trị lâu dài không? \textit{(Are you ignoring a hard task that holds long-term value?)}

2. Em nhìn cơ hội qua vẻ ngoài hay qua tiềm năng thật của nó? \textit{(Do you judge opportunities by appearance or by real potential?)}
\end{tuhocnhanh}
\ngancach